\documentclass[11pt,letterpaper]{article}
\usepackage[margin=0.88in]{geometry}
\usepackage[T1]{fontenc}
\usepackage{lmodern,microtype,amsmath,amssymb,booktabs,tabularx}
\usepackage{xcolor,graphicx,pgfplots,caption,fancyvrb}
\usepackage{hyperref}
\pgfplotsset{compat=1.18}
\definecolor{ink}{HTML}{193D52}
\definecolor{rust}{HTML}{AC4D24}
\hypersetup{colorlinks=true,urlcolor=ink,linkcolor=ink,citecolor=ink,
  pdftitle={The Consequences of Continuing to Borrow},
  pdfsubject={Fiscal pressure, inflation sensitivity, and the legacy of expensive debt}}
\setlength{\parindent}{0pt}
\setlength{\parskip}{6pt}
\captionsetup{font=small,labelfont=bf,hypcap=false}
\pgfplotsset{articleplot/.style={
  width=\linewidth,height=2.35in,
  axis lines=left, axis line style={black!45},
  tick label style={font=\small},label style={font=\small},
  scaled x ticks=false,xticklabel style={/pgf/number format/1000 sep={}},
  xmin=2026,xmax=2057,xtick={2030,2040,2050,2056},
  ymajorgrids=true,grid style={black!10},
  legend style={draw=none,font=\small,fill=none},
  legend cell align=left}}
\newcommand{\src}[2]{\href{#1}{#2}}
\newcommand{\pathfile}[1]{\nolinkurl{#1}}

\begin{document}
{\LARGE\bfseries\color{ink} The Consequences of Continuing to Borrow\par}
\vspace{4pt}
{\large Three findings about fiscal pressure, inflation, and the debt a crisis leaves behind\par}
{\small September 6, 2026\par}
\vspace{8pt}

\textquotedblleft This won't end well\textquotedblright{} is not an explanation.
If a government keeps borrowing to cover its spending and the interest on its
existing debt, what actually happens? Does the debt eventually become
manageable? Do prices spiral upward? Do investors demand higher rates, making
the problem worse?

The question becomes more interesting when we take politics seriously.
Assume that, before sustained economic pain changes the political calculation,
elected officials keep financing scheduled commitments rather than enact a
major deficit-reduction package. Borrowing is the initial response. This does
not mean that taxes, benefits, or other policies can never change. It means
we should not make a crisis disappear by assuming that politicians prevent it.

A useful answer separates what follows from accounting from what depends on
economic behavior. A quarterly simulation, using U.S.\ budget projections and
Treasury debt records, produces three findings: fiscal pressure can accumulate
without a market panic; an inflation spiral is not a mechanical consequence of
higher interest payments; and expensive borrowing can outlive the conditions
that caused it.

These are conditional findings, not a forecast of a crisis date. The relevant
data, equations, assumptions, and reproduction instructions are in the appendix.
No earlier article is needed.

\section*{1. The pressure can build before anything dramatic happens}

Start with the simplest useful equation:
\[
 \text{next debt}
 =\text{current debt}+\text{primary deficit}+\text{interest}.
\]
The primary deficit is spending other than interest, minus revenue. Replacing
a maturing bond with a new bond does not, by itself, increase the debt:
one obligation replaces another. Borrowing to pay the interest or to finance
additional spending does increase it.

But a growing dollar amount is not sufficient evidence of a breakdown. The
economy can grow too. Let \(b_t\) be debt divided by annual GDP, \(d_t\) the
primary deficit divided by that year's GDP, \(i_t\) the average nominal interest
rate on opening debt, and \(g_t\) nominal GDP growth. In a simple annual model,
\begin{equation}
 b_t=\frac{1+i_t}{1+g_t}b_{t-1}+d_t.
 \label{eq:annual}
\end{equation}
\textquotedblleft Nominal\textquotedblright{} means measured in dollars,
without removing inflation.

This equation explains how continued borrowing can coexist with a stable debt
burden. With a constant 3\% interest rate, 4\% nominal growth, and a primary
deficit of 1\% of GDP, the debt ratio converges to 104\%:
\[
 \frac{1.03}{1.04}\times 1.04+0.01=1.04.
\]
At that ratio, both the debt and GDP grow 4\% a year. The government continues
borrowing, including to pay interest, but its debt does not outgrow its income
base. This is an arithmetic counterexample to inevitable collapse, not a
prediction that investors will keep lending to the United States at those rates.

Does a U.S.-based reference path show such a settling down? Not over the next
30 years in this exercise. The simulation uses the Congressional Budget
Office's February 2026 budget and economic projections, with explicit
extensions described in Appendix A. It keeps scheduled commitments financed
and adds no discretionary fiscal rescue.
The inputs are a fixed projection vintage, not a continuously updated forecast.
\src{https://www.cbo.gov/publication/62105}{[CBO, ten-year outlook]}
\src{https://www.cbo.gov/publication/62044}{[CBO, long-term data]}

Without an added rate shock, modeled debt held by the public rises from an
opening 99.22\% of GDP to 172.77\% in the third quarter of 2056.
It is still rising at the end.

\begin{center}
\begin{minipage}{0.94\linewidth}
\begin{tikzpicture}
\begin{axis}[articleplot,ylabel={Debt held by the public (\% of GDP)},
  ymin=90,ymax=185,ytick={100,120,140,160,180}]
\addplot[ink,very thick] table[x=plot_year,y=reference_debt_pct,col sep=comma]
  {../data/processed/borrowing_article_plot_data.csv};
\end{axis}
\end{tikzpicture}
\captionof{figure}{Fiscal pressure without an added shock.}
{\footnotesize Own simulation, 2026Q4--2056Q3. End-quarter debt is divided by
that quarter's annual-rate GDP. This is a reference projection, not a modeled
test of investors' willingness to finance it.}
\end{minipage}
\end{center}

The final year's arithmetic makes the mechanism visible. Across its last four
quarters, primary deficits add 2.20 percentage points to the debt ratio;
modeled interest adds 6.88 points; and nominal GDP growth removes 5.95 points.
The net increase is 3.13 percentage points. Growth is doing substantial work.
It is simply not doing enough to offset the new obligations.
These are contributions to the change in the debt ratio, not three conventional
annual budget shares; Appendix B gives the exact calculation.

There is an important limit to this result. The reference path supplies
borrowing rates and economic growth; it does not ask investors whether they
will accept every future auction. It therefore cannot demonstrate that a
shock-free path is economically feasible indefinitely. Nor does its rising
endpoint prove that stabilization after 2056 is impossible.

What it does establish is concrete: under these financing and growth
assumptions, continued borrowing does not produce a stable debt ratio within
the projection. No initiating panic is needed for the burden to grow.

\section*{2. More interest income does not automatically mean an inflation spiral}

Interest paid by the government becomes income for someone else. If recipients
spend part of it, that spending can support demand. This suggests a troubling
feedback: higher rates increase interest payments, those payments support
spending and inflation, and the central bank raises rates again.

But the existence of that channel does not determine its net effect. Some
interest is saved or received abroad. Higher rates can also restrain borrowing,
reduce spending, and lower the market value of existing bonds. The Bank for
International Settlements discusses these competing income, credit, and
valuation channels; it does not treat interest income as a sufficient
explanation of the overall inflation response.
\src{https://www.bis.org/publications/aer-2026/high-public-debt-shifting-financial-markets}
{[BIS, 2026, chapter II]}

To see why the distinction matters, consider two deliberately severe stress
scenarios. Both add five percentage points to Treasury issuance rates for ten
years, beginning in 2026Q4. Both also add two percentage points to private
credit spreads for the first two years. Both let some additional interest
income feed into spending, let the central bank respond to inflation and
output, and continue financing the primary deficit. Neither adds a fiscal
consolidation package.

Change just one coefficient: how much of an inflation deviation survives into
the next quarter.

\begin{center}
\begin{minipage}{0.94\linewidth}
\begin{tikzpicture}
\begin{axis}[articleplot,ylabel={Annualized inflation (\%)},
  ymin=1,ymax=9.2,ytick={2,4,6,8},
  legend pos=north west]
\addplot[rust,very thick] table[x=plot_year,y=persistent_inflation_pct,col sep=comma]
  {../data/processed/borrowing_article_plot_data.csv};
\addlegendentry{98\% quarterly persistence}
\addplot[ink,very thick,dashed] table[x=plot_year,y=faster_fading_inflation_pct,col sep=comma]
  {../data/processed/borrowing_article_plot_data.csv};
\addlegendentry{90\% quarterly persistence}
\end{axis}
\end{tikzpicture}
\captionof{figure}{The same shocks, very different inflation paths.}
{\footnotesize Own simulations. Inflation is annualized quarterly GDP-price
inflation, not twelve-month consumer-price inflation. All parameters except
inflation persistence are identical. The Treasury premium ends in 2036Q3;
the private credit-spread shock ends in 2028Q3.}
\end{minipage}
\end{center}

\begin{center}
\begin{tabular}{lrr}
\toprule
 & \multicolumn{2}{c}{2056Q3 result}\\
\cmidrule(lr){2-3}
Quarterly inflation persistence & Inflation & Debt/GDP\\
\midrule
98\% & 8.04\% & 255.61\%\\
90\% & 2.89\% & 259.92\%\\
\bottomrule
\end{tabular}
\end{center}

The second case does not merely postpone a large inflation spike: inflation
never exceeds 3.15\% during the simulated period. Yet debt ends slightly
higher relative to GDP than in the high-inflation case. In these experiments,
avoiding sustained high inflation does not resolve the fiscal problem.

The apparent smallness of the parameter change is misleading. In isolation,
98\% quarterly persistence gives an inflation deviation a half-life of about
8.6 years; 90\% gives it a half-life of about 1.6 years. These are substantially
different assumptions about how the economy works. Neither coefficient has
been estimated here, and the two paths are not a confidence interval.

There is another limitation worth stating plainly. In these runs, the
Treasury-specific premium raises federal financing costs but does not directly
tighten private credit. Private restraint comes from the separate credit-spread
shock and the policy-rate response. Bond-price losses and several other
channels are also absent. Consequently, the high-inflation path is a
conditional mechanism illustration, not evidence that higher Treasury yields
would have that net effect in the actual economy.

The lesson is not that inflation is harmless, or that debt is merely a number
on paper. It is that paying interest is insufficient to derive an inflation
spiral. The fiscal ledger tells us the payment. Economic assumptions determine
what happens after the payment reaches its recipient.

\section*{3. A recovery inherits the debt issued during the stress}

Suppose sustained inflation, a recession, or financing stress eventually
changes what voters and politicians will accept. Policy changes become
possible. Market confidence improves, and borrowing rates fall.

That improvement matters immediately for new borrowing. It does not rewrite
the coupons on outstanding fixed-rate bonds. Treasury bonds pay their fixed
rate until maturity.
\src{https://www.treasurydirect.gov/marketable-securities/treasury-bonds/}
{[TreasuryDirect, Treasury bonds]}

The simplest example needs no macroeconomic simulation:
\[
 \$1\text{ trillion of fixed-rate debt at }8\%
 \quad\longrightarrow\quad \$80\text{ billion of annual interest}.
\]
If the rate available on new bonds falls to 4\%, the existing obligation still
pays \$80 billion a year while it remains outstanding. New debt of the same
face value issued at par would pay \$40 billion. The Treasury cannot simply
substitute the cheaper coupon. Buying back the old bond would require paying
its market price, which reflects its valuable remaining payments.

Actual federal debt mixes short-term bills, longer fixed-rate securities,
floating-rate notes, and inflation-protected debt. Some costs adjust quickly;
others persist. The simulation tracks those different instruments precisely
because today's market yield and the average cost of the debt stock are not
the same number. Appendix B describes that machinery. The fixed-coupon example
establishes the carryover mechanism without assuming any particular recovery
date or policy package.

Continuing to borrow during stress can leave two legacies: more principal and
high contractual payments on debt issued while rates were elevated. If those
payments also have to be borrowed, they add further debt. Lower rates help,
but recovery starts from that inherited balance sheet, not from the balance
sheet the government had before the stress.

This is not a claim that the cost is permanent or that every maturity strategy
has the same outcome. Short maturities pass higher rates through quickly and
also benefit sooner when rates fall. Long fixed-rate borrowing delays repricing
but locks in the rate available when the bond is issued. The timing is part
of the problem, not a detail that can be replaced with a single market yield.

\section*{What we can say without pretending to know the ending}

The relevant question is not whether the United States must run out of
dollars on a calculable date. It is what happens to the debt burden, prices,
financing conditions, and political choices while borrowing continues.

This exercise gives a more specific answer than a warning that things will
end badly. Under the reference assumptions, the debt burden keeps increasing
even without an added shock. Under the specified rate stress, that burden
becomes much larger whether inflation remains high or subsides. And if severe
pain eventually prompts a response, the government still has to deal with the
obligations accumulated before that response.

It does not identify the probability of the stress, a debt ratio that triggers
a run, or the point at which politics changes. The unadjusted paths describe
continued borrowing, not a prediction that decades of distress would leave
policy unchanged. They show consequences that a subsequent response would
have to confront.

The debt arithmetic is knowable. The ending depends on how the economy and
politics respond. Useful scenarios make that dependence explicit, replacing
a vague warning with an explanation.

\clearpage
\appendix
{\Large\bfseries\color{ink} Appendix: the calculations behind the article\par}

This appendix specifies the three runs used above. All reported scenario
results are calculations from the accompanying simulator, not results claimed
by CBO, Treasury, or BIS. No separate crisis narrative, fiscal-reform scenario,
or previous paper is required to interpret them.

\section{Inputs, scope, and the three experiments}

\textbf{Dates and units.}
The simulation runs quarterly from 2026Q4 through 2056Q3, inclusive.
Stocks and quarterly flows are in billions of nominal dollars. GDP is a
seasonally adjusted annual rate (SAAR). Rates, output gaps, and GDP shares are
decimals in the equations. Inflation is the annualized quarterly change in
the GDP price level; TIPS use a separate CPI-based indexation input.

\textbf{Starting balance sheet.}
Treasury's September 30, 2025 Monthly Statement of the Public Debt supplies
security-level maturities and instrument types. Marketable securities are
scaled within each class to the statement's public holdings; auction records
supply available effective-rate information, with coupon-rate fallback.
Other publicly held debt is retained as a coarse aggregate. The stock is
advanced under the reference assumptions through 2026Q3 before the experiments
start. Opening debt is \$32,215.665 billion and opening annual-rate GDP is
\$32,468.269 billion, giving 99.221998\%.
\src{https://fiscaldata.treasury.gov/datasets/monthly-statement-public-debt/}
{[Treasury, debt data]}
\src{https://fiscaldata.treasury.gov/datasets/treasury-securities-auctions-data/}
{[Treasury, auction data]}

\textbf{Budget and economy.}
The February 11, 2026 CBO vintage supplies the ten-year fiscal and quarterly
economic inputs. The separate February 25 long-term release extends the
budget and economy through 2056. CBO's baseline convention continues scheduled
benefits beyond trust-fund exhaustion; its payable-benefits alternative
differs. Thus the article studies financing scheduled commitments, not literal
legislative inaction under every existing financing restriction.
\src{https://www.cbo.gov/publication/62105}{[CBO, appendices B and C]}
\src{https://www.cbo.gov/publication/62044}{[CBO, long-term release]}

Annual primary deficits are split into four equal nominal fiscal quarters.
For post-2036 annual economic levels, constant within-year quarterly growth is
chosen so the mean of the four levels matches the supplied annual level.
This is done for nominal and real GDP and the GDP and CPI price levels.
Quarterly Treasury issuance-rate inputs after 2036 are held at the final
ten-year values. Other financing adjustments are set to zero after FY2036
because the long-term input does not supply them. These two extensions are
modeling choices, not additional CBO forecasts. The February 25 release
excludes the effects of the February 20 tariff ruling; this exercise does not
refresh that vintage.
\src{https://www.cbo.gov/publication/62044}{[CBO, release scope]}

\begin{center}
\small
\begin{tabularx}{\linewidth}{@{}lXXX@{}}
\toprule
Setting & Reference & Persistent inflation & Faster-fading inflation\\
\midrule
Treasury premium & None & 5 pp, first 40 quarters & Same\\
Private credit spread & None & 2 pp, first 8 quarters & Same\\
Inflation persistence \(\rho\) & 0.98 & 0.98 & 0.90\\
Discretionary primary adjustment & None & None & None\\
Debt-dependent yield feedback & Off & Off & Off\\
Issuance-capacity ceiling & Off & Off & Off\\
\bottomrule
\end{tabularx}
\end{center}

All cases use the same price-stability policy rule and issuance strategy below.
There is no imposed supply shock, default, buyback, yield cap, or incremental
central-bank purchase program. All auctions are assumed to be financed at the
specified rates. The code's numerical rate bounds never bind in these runs.
The reference-relative feedback system has zero deviations without a shock;
the no-shock case is therefore a fiscal projection, not an independently
identified equilibrium of investor demand.

\section{Debt accounting, maturities, and the no-shock calculation}

Let \(B_t\) be closing publicly held debt, \(D_t\) the quarterly primary deficit,
\(I_t\) modeled quarterly interest cost including TIPS principal compensation,
and \(A_t\) other financing adjustments. In the selected runs,
\begin{equation}
 B_t=B_{t-1}+D_t+I_t+A_t.
 \label{eq:ledger}
\end{equation}
The detailed ledger follows cohorts with principal, maturity quarter,
instrument type, and coupon or effective rate. Each quarter, it indexes TIPS,
resets floating rates, accrues interest on the inherited stock, and replaces
maturing principal. New issuance at quarter-end starts accruing interest in
the next quarter. Quarterly interest accrual approximates annual effective
interest by one quarter of its annual amount; it is not a literal schedule of
every semiannual coupon payment.

For a nominal fixed-rate cohort of principal \(Q\) and effective annual rate
\(y\), quarterly interest is \(Qy/4\). For a TIPS cohort, principal becomes
\(Q'=Q(1+\pi^{\rm TIPS}_t)^{1/4}\); the change in principal is counted as
inflation compensation and real-rate interest accrues on adjusted principal.
Principal can decline with deflation; at maturity a floor protects original
principal. Floating-rate notes reset to the short Treasury issuance rate plus
their fixed spread. Treasury's instrument descriptions provide the underlying
contractual conventions.
\src{https://www.treasurydirect.gov/marketable-securities/tips/}{[TreasuryDirect, TIPS]}
\src{https://www.treasurydirect.gov/marketable-securities/floating-rate-notes/}
{[TreasuryDirect, floating-rate notes]}

TIPS indexation increases the stock directly and is not also borrowed as a
cash coupon. Net cash financing covers the primary deficit, non-indexation
interest, and other financing needs. Gross issuance also includes refinancing
maturing principal. Other publicly held debt, initially \$582.764 billion,
has no modeled maturity schedule and accrues interest using the supplied
CBO average-debt-rate input.

\begin{center}
\small
\begin{tabular}{@{}lrrl@{}}
\toprule
New debt type & New borrowing share & Tenor (quarters) & Reference rate\\
\midrule
Bills & 21.5411\% & 1 & CBO 3-month\\
Notes & 51.8181\% & 28 & CBO 10-year\\
Bonds & 17.2861\% & 80 & CBO 10-year\\
TIPS & 7.0293\% & 40 & 1.75\% real\\
Floating-rate notes & 2.3253\% & 8 & CBO 3-month + 0.10 pp\\
\bottomrule
\end{tabular}
\end{center}

Maturing principal is refinanced within its instrument type; the shares in
the table allocate additional net borrowing. The tenors and coarse rate
mapping are approximations, not a reconstruction of every future auction.
Higher market yields affect newly issued debt and floating-rate resets, not
the fixed coupons of outstanding nominal debt.

\textbf{Exact debt-ratio decomposition.}
Let \(Y_t\) be current quarterly annual-rate GDP. Subtracting the previous
debt ratio from equation~\eqref{eq:ledger} gives
\begin{equation}
 \frac{B_t}{Y_t}-\frac{B_{t-1}}{Y_{t-1}}
 =\underbrace{\frac{D_t}{Y_t}}_{\text{primary}}
 +\underbrace{\frac{I_t}{Y_t}}_{\text{interest}}
 +\underbrace{\frac{A_t}{Y_t}}_{\text{other}}
 +\underbrace{B_{t-1}\left(\frac{1}{Y_t}-\frac{1}{Y_{t-1}}\right)}_{\text{growth}}.
 \label{eq:decomposition}
\end{equation}
Summing these contributions over 2055Q4--2056Q3 and multiplying by 100 yields:
\begin{center}
\begin{tabular}{lr}
\toprule
Contribution & Percentage points\\
\midrule
Primary deficits & \(+2.204807\)\\
Modeled interest & \(+6.875718\)\\
Other financing & \(0.000000\)\\
Nominal GDP growth & \(-5.946338\)\\
\midrule
Change in debt/GDP & \(+3.134187\)\\
\bottomrule
\end{tabular}
\end{center}

Each term uses its own quarter's GDP denominator. In particular, the interest
line is neither an annual flow divided by one annual GDP figure nor CBO's
net-interest budget measure. Interest received by the Federal Reserve is
included in the debt engine; the stress feedback is not a full consolidation
of Treasury, reserve-interest expense, and Fed remittances.

The article's 172.77\% endpoint also uses end-quarter annual-rate GDP.
Using mean GDP over FY2056 instead gives approximately 175.06\% in this
reference run. Denominator conventions explain why these two descriptions
of the same modeled closing stock differ.

\textbf{The stable-borrowing example.}
For the annual abstraction, set \(A=0\), \(I=iB_{t-1}\), \(D=dY_t\), and
\(Y_t=(1+g)Y_{t-1}\). Equation~\eqref{eq:ledger} then gives
equation~\eqref{eq:annual}. With fixed \(g>i>-1\), its stable fixed point is
\[
 b^*=\frac{d(1+g)}{g-i}.
\]
The distance from that fixed point shrinks by \((1+i)/(1+g)\) each year.
The 104\% example assumes rates and growth do not deteriorate as debt rises;
it is not an estimated financing schedule or a welfare assessment.

\section{The macroeconomic feedback used in the stress cases}

A superscript \(0\) denotes the no-shock reference path. Let \(x_t\) be the
fractional output gap relative to reference real output, \(\pi_t\) annualized
inflation, and \(e_t=\pi_t-\pi^0_t\). The experiments begin with \(x=e=0\).
The coefficients below are illustrative assumptions, not estimated causal
effects.

\subsection*{GDP and the primary deficit}

Write \(R_t\) for real GDP and \(Y_t\) for nominal GDP. If
\(r^0_t\) is the annualized reference real growth rate, then
\begin{align}
 R_t&=R_{t-1}(1+r^0_t)^{1/4}
       \frac{1+x_t}{1+x_{t-1}},\\
 Y_t&=Y_{t-1}(1+r^0_t)^{1/4}
       \frac{1+x_t}{1+x_{t-1}}(1+\pi_t)^{1/4}.
\end{align}
Let \(Y^{\rm pot}_t\) follow the same reference real growth and scenario
inflation but omit the output-gap factor. The primary deficit is
\begin{equation}
 D_t=D^0_t\frac{Y_t}{Y^0_t}-0.45\,x_t\frac{Y^{\rm pot}_t}{4}.
 \label{eq:primary}
\end{equation}
Thus the reference primary deficit scales with scenario nominal GDP, with
an automatic-stabilizer adjustment that reduces the deficit when output is
above reference and increases it when output is below. There is no
discretionary fiscal correction. This aggregate rule is not a
program-by-program model of benefit indexation, medical costs, or tax brackets
under sustained high inflation.

\subsection*{Policy rates and Treasury rates}

Let \(a_t\) be the policy-rate adjustment relative to reference and let
\(q^0_t\) denote the reference short rate:
\begin{align}
 a_t&=0.5a_{t-1}+0.5(1.5e_t+0.5x_t),\qquad a_{-1}=0,\\
 q_t&=\operatorname{clip}(q^0_t+a_t,\,0,\,0.99),\\
 u_t&=(q_t-\pi_t)-(q^0_t-\pi^0_t).
\end{align}
The demand equation uses \(u_t\) as a proxy for real monetary tightening,
based on current inflation rather than modeled inflation expectations.
Although the policy configuration labels a 2\% inflation target, the active
reaction here is to deviations from \(\pi^0_t\).

The Treasury premium \(h_t\) equals 0.05 for the first 40 stress quarters and
zero thereafter. New issuance rates are
\begin{align}
 y^{\rm short}_t&=\operatorname{clip}(q_t+h_t,\,0,\,0.99),\\
 y^{\rm note}_t&=\operatorname{clip}(y^{0,\rm note}_t+0.75a_t+h_t,\,0,\,0.99),\\
 y^{\rm bond}_t&=\operatorname{clip}(y^{0,\rm bond}_t+0.50a_t+h_t,\,0,\,0.99),\\
 y^{\rm TIPS,real}_t&=\operatorname{clip}(0.0175+0.50a_t+h_t,\,-0.99,\,0.99).
\end{align}
Floating-rate resets use \(y^{\rm short}_t\) plus the cohort's spread.
The coarse other-debt aggregate retains its supplied average rate.
For TIPS indexation, the reference lagged CPI inflation input is increased by
\(e_{t-1}\), initially zero. Using the GDP-inflation deviation as the CPI
deviation is another approximation.

\subsection*{Interest recipients and the demand impulse}

The spending fractions apply to incremental modeled cash interest, not
to all debt issuance or all accumulated wealth:
\begin{center}
\small
\begin{tabular}{lrr}
\toprule
Recipient & Allocated share & Fraction spent on domestic demand\\
\midrule
Domestic private sector & 56.8805\% & 25\%\\
Foreign holders & 28.9732\% & 5\%\\
Federal Reserve & 14.1463\% & 0\%\\
\bottomrule
\end{tabular}
\end{center}

The allocation uses pinned CBO FY2026 debt and Fed-holding inputs and
June 2026 Treasury International Capital foreign holdings. Different dates
and valuation conventions make these shares proxies, not exact holder-by-holder
cash-flow measurements. The spending fractions themselves are assumptions;
zero for the Fed suppresses its direct consumption channel, not all possible
fiscal or monetary effects of its balance sheet.
\src{https://github.com/US-CBO/cbo-data}{[CBO, official data repository]}
\src{https://ticdata.treasury.gov/resource-center/data-chart-center/tic/Documents/slt_table5.html}
{[Treasury, foreign holdings table; pinned June 2026 observation]}

The weighted domestic spending fraction is
\[
 w=(0.5688049586)(0.25)+(0.2897321139)(0.05)
   +(0.1414629275)(0)=0.1566878454.
\]
Let \(C_t\) be the modeled cash-interest proxy, excluding the TIPS principal
indexation accrual. The engine treats non-indexation effective interest cost
as a quarterly cash flow, rather than using the actual coupon calendar.
Define annualized incremental impulses
\begin{align}
 \Delta I_t&=w\left(\frac{4C_t}{Y_t}-\frac{4C^0_t}{Y^0_t}\right),\\
 \Delta P_t&=\frac{4}{Y_t}
   \left(D_t-D^0_t\frac{Y_t}{Y^0_t}\right).
\end{align}
Negative deviations are allowed. The roughly 15.67\% weight is not an estimate
that this share of all federal interest payments necessarily becomes
additional aggregate demand.

\subsection*{Output and inflation dynamics}

Let \(s_t=0.02\) for the first eight stress quarters and zero thereafter.
The next quarter's output gap and inflation deviation are
\begin{align}
 x_{t+1}
 &=0.80x_t-0.25u_{t-1}-0.15s_{t-1}
     +0.80\Delta P_t+\Delta I_t, \label{eq:output}\\
 e_{t+1}&=\rho e_t+0.06x_{t-1}. \label{eq:inflation}
\end{align}
Unavailable pre-start lags are zero. The first experiment uses \(\rho=0.98\);
the sensitivity uses \(\rho=0.90\). For an isolated deviation with no new
forcing, the half-life in years is
\[
 H(\rho)=\frac{\log(0.5)}{4\log(\rho)},
 \qquad H(0.98)=8.5774,\quad H(0.90)=1.6447.
\]
For a constant externally held output gap, equation~\eqref{eq:inflation}
would give \(e=0.06x/(1-\rho)\): \(3x\) versus \(0.6x\).
This is an explanation of the coefficient's importance, not a solution of
the full coupled system.

Within each quarter, the engine advances GDP using the current gap and
inflation, computes the primary deficit and policy response, services and
refinances the debt, computes the two fiscal impulses, and then updates
next-quarter output and inflation. A parallel reference ledger supplies
the reference interest amounts.

\textbf{What these equations omit.}
The Treasury premium does not itself appear in
equation~\eqref{eq:output}; its demand effect comes through interest income,
while the distinct private spread and monetary response restrain demand.
The system does not estimate debt-sensitive investor demand, an auction
failure threshold, expected future fiscal policy, exchange-rate responses,
private-capital accumulation, productivity damage, bondholder valuation
losses, or endogenous political change. It is not a complete equilibrium
model of household wealth and government liabilities. These omissions limit
interpretation of both inflation and real activity, especially far from the
reference path.

\section{Reproduction and verification}

The accompanying reproduction script runs only the three cases in this
article and uses existing pinned inputs. From the project root:
\begin{Verbatim}[fontsize=\small]
uv run python scripts/reproduce_borrowing_article.py
latexmk -pdf -cd -interaction=nonstopmode -halt-on-error \
  docs/continuing_to_borrow.tex
\end{Verbatim}
An already installed project environment can instead run the script with
\pathfile{.venv/bin/python}. Building the PDF also requires a LaTeX installation
with PGFPlots. No network data refresh is part of the calculation.

\textbf{Required input snapshots.}
The following files are in \pathfile{data/processed/}:
\begin{Verbatim}[fontsize=\small]
cbo_baseline_fy_2026-02.csv
cbo_economy_quarterly_2026-02.csv
cbo_long_term_budget_fy_2026-02-25.csv
cbo_long_term_economy_annual_2026-02-25.csv
treasury_securities_2025-09-30.csv
initial_conditions_2025-09-30.csv
\end{Verbatim}
The loader also reads two provenance manifests in \pathfile{data/metadata/}:
\begin{Verbatim}[fontsize=\small]
manifest_cbo-2026-02_treasury-2025-09-30.json
manifest_cbo-long-term-2026-02-25.json
\end{Verbatim}
Recipient shares come from
\pathfile{config/calibration/interest_recipient_shares_2026.json}.
The manifests identify the official source files and normalization; the
snapshots were retrieved on August 22, 2026. Source links in this article
may subsequently show newer information.

\textbf{Relevant implementation.}
The script calls the existing \pathfile{src/debt_sim/} modules:
\pathfile{data.py} and \pathfile{scenarios.py} prepare the inputs;
\pathfile{instruments.py}, \pathfile{debt_stock.py}, and \pathfile{model.py}
service and refinance the cohorts; \pathfile{crisis.py} implements the
quarterly feedback described in Appendix C; \pathfile{policy.py} supplies
the policy settings; and \pathfile{sustainability.py} supplies the exact
debt-ratio decomposition. The appendix explains the selected mechanisms;
the accompanying source and numerical snapshots supply exact replication.

\textbf{Outputs.}
All new outputs use the prefix \pathfile{data/processed/borrowing_article_}.
The files \pathfile{paths.csv}, \pathfile{summary.csv}, and
\pathfile{decomposition.csv} contain the quarterly runs, endpoint summaries,
and reference debt-ratio contributions. The figures read
\pathfile{plot_data.csv}. The exact 120-quarter reference inputs are written
to \pathfile{reference_inputs.csv}. The file \pathfile{metadata.json} records
all rule settings, shock paths, opening stocks, and SHA-256 fingerprints
of the processed input snapshots, recipient configuration, reproduction
script, and simulator Python modules.

The script checks that the debt ledger reconciles within \(10^{-7}\) billion
dollars, that the debt-ratio decomposition reconciles within \(10^{-10}\),
that no numerical rate bound binds, and that discretionary fiscal adjustments
remain zero. It verifies the following rounded endpoints rather than fitting
parameters to obtain them:
\begin{center}
\begin{tabular}{lrr}
\toprule
Run & Inflation, 2056Q3 & Debt/GDP, 2056Q3\\
\midrule
Reference & 2.06\% & 172.77\%\\
Persistent inflation & 8.04\% & 255.61\%\\
Faster-fading inflation & 2.89\% & 259.92\%\\
\bottomrule
\end{tabular}
\end{center}
These checks establish reproduction and accounting consistency, not empirical
validation of the macroeconomic coefficients. The fixed-coupon example in
Section 3 is direct multiplication and does not use an assumed recovery path.

\end{document}
